% ----------------------------------------------------------------------
%  Итоги главы 2: сводка, ответы, задачи, литература
% ----------------------------------------------------------------------
\section*{Итоги главы}
\addcontentsline{toc}{section}{Итоги главы}
\markright{Итоги главы 2}

\begin{table}[ht]
\centering
\caption{Точечные дефекты и дислокации: сравнение}
\label{tab:ch2-summary}
\small
\begin{tabularx}{\textwidth}{@{}>{\raggedright\arraybackslash}p{30mm}>{\raggedright\arraybackslash}X>{\raggedright\arraybackslash}X@{}}
\toprule
 & Точечные дефекты (вакансии) & Дислокации \\
\midrule
Энергия & $\EvF\sim1\eV$ на дефект & $\sim\frac12Gb^2$ на единицу длины, $\sim5$--$8\eV$ на атомную плоскость \\
Равновесие & есть равновесная концентрация $C_v=e^{-\EvF/kT}$, у $\Tpl$ $\sim10^{-4}$ & термодинамически неравновесны; плотность $10^{6}$--$10^{12}$~см$^{-2}$ определяется историей \\
Движение & термически активированные перескоки, $\Evm\sim0{,}5$--$1\eV$; межузельные атомы~— $\sim0{,}1\eV$ & скольжение (при любой $T$, под напряжением); переползание (диффузионное) \\
Поле напряжений & короткодействующее ($\propto r^{-3}$) & дальнодействующее ($\propto r^{-1}$) \\
Чем важны & диффузия, отжиг, радиационные повреждения, переползание дислокаций & пластичность, упрочнение, полигонизация, рост кристаллов \\
\bottomrule
\end{tabularx}
\end{table}

\begin{itogi}[Главные формулы главы]
\begin{itemize}
\item Равновесная концентрация: $C=\exp(S/k)\exp(-E^{\mathrm F}/kT)$;
$\EvF\approx10\,k\Tpl$.
\item Самодиффузия: $D=a^2\nu_0\exp\bigl(-(\EvF+\Evm)/kT\bigr)$; $D_{\text{сд}}=C_vD_v$.
\item Дивакансии: $C_{2v}=\frac z2C_{1v}^2\exp(B_2/kT)$.
\item Отжиг: $\tau=\alpha r^2/D_v$; соотношение Эйнштейна $\mu=D/kT$.
\item Теоретическая прочность $\tau_{\text{теор}}\approx G/2\pi$; напряжение
Пайерлса $\tau_\text{П}\approx(2G/K)\exp(-2\pi w/b)$.
\item Формула Орована: $\varepsilon=\rho bl$, $\dot\varepsilon=\rho bv$.
\item Поля: $\sigma_{\theta z}=Gb/2\pi r$ (винтовая); $D=Gb/2\pi(1-\nu)$ (краевая).
\item Энергия дислокации $E\approx\frac{Gb^2}{4\pi K}\ln\frac R{r_0}\approx\alpha Gb^2$.
\item Сила Пича--Кёлера $\vb f=(\tn\sigma\cdot\vb b)\times\vb t$; сила
скольжения $\tau b$.
\item Атмосферы: $C=C_0\exp(E/kT)$, $T_\text{н}=E_{\max}/k\ln(1/C_0)$.
\item Источник Франка--Рида: $\tau_{\text{кр}}\approx Gb/l$.
\item Скопление: концентрация напряжений в $N$ раз; Холл--Петч
$\tau_\text{т}=\tau_0+kl^{-1/2}$. Стенка: $\varphi=b/h$.
\item Критерий Франка: $\sum b^2_{\text{до}}>\sum b^2_{\text{после}}$;
ширина расщепления $d_0\propto Gb^2/\gamma$.
\end{itemize}
\end{itogi}

%=======================================================================
\section*{Ответы на вопросы для самопроверки}
\addcontentsline{toc}{section}{Ответы на вопросы для самопроверки}
\markright{Ответы на вопросы главы 2}

{\small
\paragraph{2.1.2}
\begin{enumerate}[label=\arabic*.]
\item Не учтены релаксация соседних атомов к вакансии и перераспределение
электронов; оба эффекта снижают энергию в 2,5--3 раза.
\item $C(1300)/C(1000)=\exp\bigl[\frac{\EvF}{k}\bigl(\frac1{1000}-\frac1{1300}\bigr)\bigr]
=\exp(14\,850\cdot2{,}31\cdot10^{-4})\approx e^{3{,}4}\approx30$ раз
(для $\EvF=1{,}28\eV$).
\item $\EiF$ в 3--4 раза больше $\EvF$, и больцмановский множитель для
межузельных атомов на много порядков меньше (у точки плавления меди
$\sim10^{-11}$).
\item Перескок атома медленнее в $1/C_v$ раз: ему нужно ещё дождаться
вакансии в соседнем узле. У точки плавления~— в $\sim10^4$ раз, при более
низких $T$~— намного больше.
\item Атом перескакивает, только если рядом вакансия (вероятность
$\propto e^{-\EvF/kT}$), и только преодолев барьер ($\propto
e^{-\Evm/kT}$); вероятности перемножаются, энергии складываются.
\item Эстафетный механизм с промежуточным минимумом энергии даёт очень
низкий барьер: $\Eim\sim0{,}1\eV$.
\end{enumerate}

\paragraph{2.1.3}
\begin{enumerate}[label=\arabic*.]
\item При удалении атома, соседнего с вакансией, рвётся на одну связь
меньше; выигрыш~— энергия связи $B_2$.
\item Дивакансия~— «встреча» двух вакансий; вероятность встречи
пропорциональна произведению их концентраций.
\item У дивакансии в седловой точке меньше атомов, создающих отталкивание,
и атом может пройти по более свободному пути. Тетраэдрическая
тривакансия для перемещения должна частично «разобраться», что требует
больших смещений атомов ($\approx1{,}9\eV$).
\end{enumerate}

\paragraph{2.1.4}
\begin{enumerate}[label=\arabic*.]
\item Нужна энергия $\EvF+\EiF\approx4$--$5\eV$; больцмановский множитель
ничтожен.
\item $\tau\propto r^2\propto1/\rho$: время уменьшится примерно в 100 раз
(с небольшой поправкой на логарифм в $\alpha$).
\item Сравниваются состояния с одинаковой концентрацией, и множитель
$F(C)$ сокращается.
\item Энергия активации зависит от $C$: одновременно идут несколько
процессов с разными энергиями активации.
\end{enumerate}

\paragraph{2.1.5}
\begin{enumerate}[label=\arabic*.]
\item Даже у точки плавления их равновесная концентрация ничтожна~—
замораживать нечего.
\item Стоков (границ) становится больше, $n$ уменьшается, $\gamma=\beta n/\nu_0$
падает~— закалка менее эффективна.
\item $T_{\max}=\frac{4\cdot1\cdot56}{57^2}\cdot1~\text{МэВ}\approx0{,}069$~МэВ
$\approx69$~кэВ.
\item Межузельные атомы переносятся цепочками замещающих соударений вдоль
плотноупакованных рядов; у вакансий такого механизма нет.
\end{enumerate}

\paragraph{2.1.6}
\begin{enumerate}[label=\arabic*.]
\item Каждая вакансия добавляет узел (атомный объём), а параметр решётки
меняется лишь из-за малой релаксации.
\item Отрицательный: $\Delta L/L-\Delta a/a=-C/3$.
\item При высоких температурах преобладает тепловая (фононная) часть
сопротивления; при низких остаётся остаточная, её можно измерить точно.
\item Во время закалки часть вакансий уходит на стоки и объединяется в
дивакансии и комплексы.
\end{enumerate}

\paragraph{2.2.1}
\begin{enumerate}[label=\arabic*.]
\item Сдвиг идёт не одновременно по всей плоскости, а последовательно~—
пробегом дислокаций, что требует малых напряжений.
\item $\vb b\cdot\vb t=0$~— краевая; $\vb b\parallel\vb t$~— винтовая;
иначе смешанная.
\item Да: вектор Бюргерса постоянен, а ориентация линии меняется
(например, у петли).
\item $\vb b$ постоянен вдоль линии; контур Бюргерса вокруг «конца» можно
было бы стянуть в точку, не пересекая дефекта, и он должен был бы
замкнуться~— противоречие.
\item $\vb b_3=-(\vb b_1+\vb b_2)$.
\end{enumerate}

\paragraph{2.2.2}
\begin{enumerate}[label=\arabic*.]
\item Скольжение не требует переноса массы; переползание требует
диффузии вакансий, которая при низких $T$ практически остановлена.
\item У плотноупакованных плоскостей наибольшее межплоскостное
расстояние $a$ и наименьший $b$, а $\tau_\text{П}\propto\exp(-2\pi a/Kb)$.
\item У винтовой дислокации $\vb b\parallel\vb t$, любая плоскость,
содержащая линию, содержит $\vb b$, и $\vb b\cdot\vb n=0$.
\item Дислокация движется перпендикулярно $\vb b$, атомы ядра смещаются
параллельно $\vb b$.
\item $v=\dot\varepsilon/(\rho b)=10^{-3}/(10^8\cdot2{,}5\cdot10^{-8})=4\cdot10^{-4}$~см/с
$=4$~мкм/с.
\end{enumerate}

\paragraph{2.2.3}
\begin{enumerate}[label=\arabic*.]
\item Поле винтовой дислокации~— чистый сдвиг без гидростатического
давления, а «сферический» дефект взаимодействует только с давлением.
\item $1/(1-\nu)=1{,}5$ раза.
\item Энергия $\sim Gb^2$ на единицу длины несопоставимо больше $kT$, и
конфигурационная энтропия не может её скомпенсировать.
\item Скольжение~— касательное напряжение в плоскости скольжения вдоль
$\vb b$; переползание~— нормальное напряжение вдоль $\vb b$.
\item Для одноимённых дислокаций положение $x=0$~— устойчивое равновесие
(сила $f_x$ возвращающая).
\item $T_\text{н}=0{,}2/(8{,}6\cdot10^{-5}\cdot\ln10^3)\approx340$~К.
\end{enumerate}

\paragraph{2.2.4}
\begin{enumerate}[label=\arabic*.]
\item Ступенька не исчезает, и не нужно зарождать двумерные островки
нового слоя.
\item Больше упругой энергии несоответствия нужно снять; расстояние между
дислокациями $\sim b/\delta$, где $\delta$~— несоответствие параметров.
\item У отрезка длиной $l$ наименьший возможный радиус кривизны $l/2$;
после этого расширение требует всё меньшего напряжения~— неустойчивость.
\item $\tau_{\text{кр}}\approx Gb/l=48\cdot10^9\cdot0{,}26\cdot10^{-9}/10^{-6}\approx12$~МПа.
\item Обратное напряжение от скопления испущенных петель у препятствия.
\end{enumerate}

\paragraph{2.2.5}
\begin{enumerate}[label=\arabic*.]
\item Его энергия больше энергии разнесённых дислокаций, а любая
перегруппировка понижает энергию; скопление создаётся только источником,
«прижимающим» дислокации к препятствию.
\item В 50 раз.
\item В мелком зерне короче скопления и меньше концентрация напряжений у
границ, поэтому для передачи деформации нужно большее приложенное
напряжение (Холл--Петч).
\item $\varphi=b/h=0{,}01$~рад~$\approx0{,}6^\circ$.
\item Нужно переползание, т.\,е. диффузия; она идёт при
$T\gtrsim(0{,}3\text{--}0{,}5)\Tpl$.
\end{enumerate}

\paragraph{2.2.6}
\begin{enumerate}[label=\arabic*.]
\item Энергия дислокации пропорциональна $b^2$.
\item Сумма векторов: $\frac a6[2+1,\,1+2,\,1-1]=\frac a2[110]$; квадраты:
$a^2/2>a^2/6+a^2/6$~— реакция выгодна.
\item Вычитание~— не хватает части слоя (одна «гексагональная» тройка);
внедрение~— лишний слой (две тройки подряд).
\item Её вектор Бюргерса перпендикулярен плоскости дефекта упаковки и не
лежит в плоскости, где лежит линия.
\item В меди: у неё $\gamma$ меньше, дислокации расщеплены шире, и
частичным дислокациям труднее слиться для поперечного скольжения.
\end{enumerate}

\paragraph{2.2.7}
\begin{enumerate}[label=\arabic*.]
\item $AC=A\delta+\delta C$: $\frac a2[10\br1]=\frac a6[2\br1\br1]+\frac a6[11\br2]$.
\item Вершинная дислокация $\frac a6\nfam{110}$ имеет вектор Бюргерса, не
лежащий ни в одной из плоскостей $\plfam{111}$, содержащих её линию, и
связана с дефектами упаковки в двух плоскостях.
\item $3a^2/4+3a^2/4=3a^2/2>a^2$~— выгодна.
\item $c$- и $(c+a)$-дислокации (только у них есть составляющая вдоль $c$).
\end{enumerate}
}

%=======================================================================
\section*{Задачи}
\addcontentsline{toc}{section}{Задачи}
\markright{Задачи к главе 2}

\begin{enumerate}[label={\sffamily\bfseries\color{fmHead}\arabic*.}]
\item Оцените равновесную концентрацию вакансий в алюминии при 600\,°C,
если $\EvF=0{,}67\eV$, $S_v=0{,}7\,k$.\\
\emph{Ответ:} $C_v=e^{0{,}7}\exp\bigl(-0{,}67/(8{,}6\cdot10^{-5}\cdot873)\bigr)\approx3\cdot10^{-4}$.
\item Коэффициент самодиффузии меди: $D_0=0{,}2$~см$^2$/с, $Q=2{,}04\eV$.
Найдите $D$ при 1000~К и характерный путь диффузии $\sqrt{Dt}$ за 1~ч.\\
\emph{Ответ:} $D\approx1\cdot10^{-11}$~см$^2$/с; $\sqrt{Dt}\approx2$~мкм.
\item Какова пороговая энергия нейтрона для смещения атома железа ($M=56$,
$E_d=25\eV$)?\\
\emph{Ответ:} $T_{\max}/E=4\cdot56/57^2\approx0{,}069$, $E_{\min}\approx360\eV$.
\item Найдите энергию единицы длины краевой дислокации в железе ($G=80$~ГПа,
$b=0{,}248$~нм, $\nu=0{,}29$, $\ln(R/r_0)=7$) и энергию на одну атомную
плоскость.\\
\emph{Ответ:} $E\approx4\cdot10^{-9}$~Дж/м; $Eb\approx6\eV$.
\item Источник Франка--Рида в алюминии ($G=26$~ГПа, $b=0{,}286$~нм) имеет
длину $0{,}5$~мкм. При каком напряжении он начнёт работать?\\
\emph{Ответ:} $\tau_{\text{кр}}\approx Gb/l\approx15$~МПа.
\item Предел текучести стали 200~МПа при размере зерна 64~мкм и 300~МПа при
16~мкм. Найдите $\sigma_0$, $k$ и предел текучести при размере зерна
4~мкм.\\
\emph{Ответ:} $k=0{,}8$~МПа$\cdot$м$^{1/2}$, $\sigma_0=100$~МПа,
$\sigma_\text{т}(4~\text{мкм})=500$~МПа.
\item Малоугловая граница наклона в меди имеет разориентацию $1^\circ$.
Найдите расстояние между дислокациями ($b=0{,}256$~нм).\\
\emph{Ответ:} $h=b/\varphi=0{,}256/0{,}01745\approx15$~нм.
\item Вакансии в зерне размером 20~мкм уходят на границы. Во сколько раз
уменьшится время отжига при повышении температуры с 300 до 350~К, если
$\Evm=0{,}7\eV$?\\
\emph{Ответ:} в $\exp\bigl[\frac{0{,}7}{8{,}6\cdot10^{-5}}\bigl(\frac1{300}-\frac1{350}\bigr)\bigr]\approx
e^{3{,}9}\approx50$ раз.
\end{enumerate}

%=======================================================================
\section*{Литература}
\addcontentsline{toc}{section}{Литература к главе 2}
\markright{Литература к главе 2}

{\small
\paragraph{Использованная в исходном учебнике}
\begin{enumerate}[label=\arabic*.]
\item Новиков И.\,И., Розин К.\,М. Кристаллография и дефекты
кристаллической решётки.~— М.: Металлургия, 1990.~— 231~с.
\item Штремель М.\,А. Прочность сплавов. Ч.~1. Дефекты решётки.~— М.:
Металлургия, 1982.~— 277~с.
\item Ван Бюрен. Дефекты в кристаллах.~— М.: Изд-во иностр. лит., 1962.~— 346~с.
\item Вопросы теории дефектов в кристаллах: сб. науч. тр.~— М.: Наука, 1987.~— 176~с.
\item Трушин Ю.\,В. Физическое материаловедение.~— СПб.: Наука, 2000.~— 286~с.
\item Дамаск А., Динс Дж. Точечные дефекты в металлах.~— М.: Мир, 1966.~— 291~с.
\item Келли А., Гровс Г. Кристаллография и дефекты в кристаллах.~— М.:
Мир, 1974.~— 496~с.
\item Фридель Ж. Дислокации.~— М.: Мир, 1967.~— 627~с.
\item Орлов А.\,Н. Введение в теорию дефектов в кристаллах.~— М.: Высшая
школа, 1983.~— 144~с.
\item Коттрелл А. Теория дислокаций.~— М.: Мир, 1969.~— 95~с.
\end{enumerate}

\paragraph{Дополнительно}
\begin{enumerate}[label=\arabic*., resume]
\item Хирт Дж., Лоте И. Теория дислокаций.~— М.: Атомиздат, 1972.
\item Hull D., Bacon D.\,J. Introduction to Dislocations. 5th ed.~—
Oxford: Butterworth-Heinemann, 2011.
\end{enumerate}
}
